\section{The geometric endpoint wall}\label{sec:geometric-wall}
We now complete a second wall mechanism in the same rank-deficient cubic
experiment. Fix $D>r>0$, strict weights $\alpha,\gamma=1-\alpha$,
positive $U=(u,u)$ and invertible positive $V$, and the exterior design.
Put
\begin{equation}\label{eq:geometric-centre}
 s_0=\frac{2rD}{3D-r},\qquad s_\delta=s_0-\delta,
 \qquad f=z(z-r)^2,\qquad g_\delta=(z-s_\delta)^3.
\end{equation}
Then $2r/3<s_0<r$. Assume the strict weight margin
\begin{equation}\label{eq:geometric-weight-margin}
 E_0=\gamma A(s_0)-4\pi r^3>0,
 \qquad A(s)=(r+3s)(2r-3s)^2.
\end{equation}
Let $\theta_\delta$ be this base centre and $P_\delta$ its observation.
Here $B(\theta_\delta)=-(3D-r)\delta$. Thus positive $\delta$ is the
identified side, while the weight constraint stays inactive locally.
Write $c_*=4r^3/A(s_0)$ and
\[
 a=\frac{4rs_0}{r+3s_0},\qquad
 h_*=(1-c_*)f+c_*g_0=(z-a)(z-D)^2,\qquad
 \beta=\gamma/c_*.
\]
The remote weight on $h_*$ is $\beta>\pi$, and its other weight is
$1-\beta>\pi$. Its first channel is $(u,u)$ and its second channel is
\begin{equation}\label{eq:geometric-remote-channel}
 v_f^*=\frac{\alpha v_1+(\gamma-\beta)v_2}{1-\beta},\qquad
 v_h^*=v_2.
\end{equation}
Both columns are positive and distinct. Denote this remote point by
$\xi_B$. Theorem~\ref{thm:cubic-fibre} gives exactly the base and this
remote representative at $\delta=0$.

\begin{proposition}[The endpoint-root parametrization]\label{prop:geometric-cone}
A full relative neighbourhood of $\xi_B$, modulo the component
permutation, has the following exact feasible parametrization:
\begin{align}
 f_v(z)&=(z-a_0)\bigl((z-r-b_0/2)^2-X\bigr),\nonumber\\
 h_v(z)&=(z-a-c_0)(z-D-d_1)(z-D-d_2),\label{eq:geometric-chart}\\
 K_B&=\{a_0,X\ge0,\quad d_1\le d_2\le0\}.\nonumber
\end{align}
The weight on $h$ is $\beta+k$, and the four binary channel coordinates
have their unrestricted small increments. The other coordinates
$b_0,c_0,k$ are free. This is a polynomial parametrization with continuous
semialgebraic inverse; no differentiability of its inverse at
$d_1=d_2$ is asserted. Its first polynomial jets are
\begin{equation}\label{eq:geometric-jets}
 p_f=-a_0(z-r)^2-b_0z(z-r)-Xz,\qquad
 p_h=-c_0(z-D)^2-(d_1+d_2)(z-a)(z-D).
\end{equation}
Let $S_B$ be the score at $\xi_B$ applied to these jets and the stochastic
increments. Then $S_B$ has no nonzero zero on $K_B$. If
$B_B=\left.\frac{d}{d\delta}P_\delta\right|_0$, then
$B_B\notin S_BK_B$. In particular
\begin{equation}\label{eq:mu-B}
 \mu_B=\frac12\min_{v\in K_B}\|S_Bv-B_B\|_{I,0}>0.
\end{equation}
The minimum is attained. After seven free coordinates are eliminated,
it is a quadratic program with three nonnegative coordinates and eight
independent-support tests.
\end{proposition}
\begin{proof}
Separated simple and double factors give the first component's simple
root, quadratic mean and squared half-separation as analytic functions
of its coefficients. At the second component the separated simple root
is analytic; the two roots near $D$ have a unique increasing ordering
and depend continuously and semialgebraically on the coefficients.
Their feasibility is exactly $d_1\le d_2\le0$. All other root bounds,
weights and channel inequalities have strict margins, except $a_0\ge0$
and $X\ge0$. This proves exact coverage and injectivity of the
parametrization, not merely containment of tangent directions. Its
polynomial formulas give a uniform $O(\tau^2)$ coefficient remainder
on substituting $\tau v$ for bounded $v\in K_B$. Positive normalizers
and probabilities give the same remainder for the observation embedding.

We prove the two score assertions directly. Put $\dot g=3(z-s_0)^2$,
the derivative of $g_\delta$ at zero. Equality of remote and base scores
implies equality of their derivatives of $q,K$ by the differentiated
normalizer argument in Lemma~\ref{lem:normalizer}. The coefficient
second marginal at the remote point is
$(1-\beta)v_f^* f+\beta v_2 h_*$. Its two coefficient columns are
independent. Modulo $\operatorname{span}(f,g_0)$, equality of its
derivative with $\gamma v_2\dot g$ gives
\[
 p_f=t_f(g_0-f),\qquad p_h=c_*\dot g+t_h(g_0-f).
\]
The polynomials have the stated degrees, so no additional monic direction
occurs. Evaluation of the first identity at $0$ and $r$ gives
$t_f=a_0r^2/s_0^3\ge0$ and $t_f=-Xr/(r-s_0)^3\le0$.
Thus $t_f=0$ and $p_f=0$.

To interpret the second identity, consider the algebraic double-root
pencil curve, without imposing its endpoint constraint,
\[
 \bar h_\delta=(1-c_*(s_\delta))f+c_*(s_\delta)g_\delta
              =(z-a(s_\delta))(z-b(s_\delta))^2,
 \qquad b(s)=\frac{rs}{3s-2r}.
\]
At zero, $b(s_0)=D$ and
\begin{equation}\label{eq:outward-speed}
 \dot b=\left.\frac{d}{d\delta}b(s_0-\delta)\right|_0
               =\frac{2r^2}{(3s_0-2r)^2}>0.
\end{equation}
Both $p_h(D)$ and $\dot{\bar h}(D)$ vanish. Since
$g_0(D)-f(D)\ne0$, this determines $t_h$ uniquely and implies
$p_h=\dot{\bar h}$. Differentiation in $z$ at $D$ now gives
\[
 -(d_1+d_2)(D-a)=p_h'(D)=-2(D-a)\dot b.
\]
This would require $d_1+d_2=2\dot b>0$, contrary to feasibility.
Therefore $S_Bv=B_B$ is impossible.

For a zero score the same second-marginal argument instead gives
$p_f=t_f(g_0-f)$ and $p_h=t_h(g_0-f)$. The first polynomial again
vanishes. Evaluating the second at $D$ forces $t_h=0$. Its coefficients
then give $c_0=0$ and $d_1+d_2=0$, whence $d_1=d_2=0$ on $K_B$.
The first polynomial's independent columns give $a_0=b_0=X=0$.
Coefficient functionals dual to $f,h_*$ recover zero derivatives of both
coefficient matrices, and their entry, row and column sums give zero
weight and channel increments. Thus the entire cone vector is zero.
Compactness of the feasible unit sphere yields
$\|S_Bv\|_{I,0}\ge c\|v\|$ on $K_B$. The objective in
\eqref{eq:mu-B} is consequently coercive and has a positive attained
minimum.

Finally the score depends on $d_1,d_2$ only through
$w=-(d_1+d_2)\ge0$. Every $w\ge0$ is realized, for example by
$d_1=d_2=-w/2$. The remaining nonnegative coordinates are $a_0,X$;
there are seven free coordinates. Their columns are independent by the
zero-score argument. Fisher-orthogonal projection off those columns
reduces the problem to three nonnegative variables. The minimal-positive-
support argument of Proposition~\ref{prop:cubic-Q} gives the eight
support tests, including the empty support and only inverting positive
 definite supported Gram matrices. This remains valid when the full
projected Gram matrix is singular.
\end{proof}

\begin{theorem}[The geometric-wall two-scale law]\label{thm:geometric-wall}
Under \eqref{eq:geometric-centre}--\eqref{eq:geometric-weight-margin},
the remote entrance satisfies
\begin{equation}\label{eq:geometric-entrance}
 d_{{\rm rem},B}(\delta)=\mu_B\delta+k_B\delta^2+o(\delta^2),
 \qquad\mu_B>0,
\end{equation}
where $k_B$ is given by Theorem~\ref{thm:critical-window} in the
parametrization \eqref{eq:geometric-chart}. In particular, for fixed
$\lambda>0$,
\begin{align}
 \lambda<\mu_B:\quad
 \omega_{\theta_\delta}(\lambda\delta)
 &=C_{\loc}(\theta_0)\sqrt{\lambda\delta}+O_\lambda(\delta),
                                                        \label{eq:B-below}\\
 \lambda>\mu_B:\quad
 \omega_{\theta_\delta}(\lambda\delta)
 &=D-r+O_\lambda(\sqrt\delta).                         \label{eq:B-above}
\end{align}
The base local Fisher constant stays positive and finite at the wall.
The strict cases in the critical window
$t=\mu_B\delta+s\delta^2$ are decided by $s-k_B$.

On the nonidentified side $\theta_{-\varepsilon}$, for all sufficiently
small $\varepsilon>0$, the additional exact pencil interval has length
\begin{equation}\label{eq:B-pencil-opening}
 c_D(s_0+\varepsilon)-c_*(s_0+\varepsilon)
       =\Lambda_B\varepsilon^2+O(\varepsilon^3),\qquad
 \Lambda_B=\frac{-R_{s_0}''(D)}{2(R_{s_0}(D)-1)^2}\dot b^{\,2}>0,
\end{equation}
where $R_s(z)=f(z)/(z-s)^3$. Nevertheless its exact whole-model
spectral diameter is constant:
\begin{equation}\label{eq:B-exact-diameter}
 \omega_{\theta_{-\varepsilon}}(0)=D-r.
\end{equation}
Thus the two wall mechanisms differ both in the size of the newborn
pencil interval and in the opening of its exact spectral diameter.
\end{theorem}
\begin{proof}
The endpoint parametrization and its score separation satisfy the
hypotheses of Theorems~\ref{thm:remote-general} and
\ref{thm:critical-window}. The derivative of the square-root observation
embedding is the score multiplied entrywise by
$\sqrt{w_j}/(2\sqrt{P_{0,je}})$, so its entrance minimum is precisely
\eqref{eq:mu-B}. This proves \eqref{eq:geometric-entrance} and the
critical-window assertion, including possible nonuniqueness of the
minimizing endpoint-root split.

The sorted spectra at the wall are $(0,s_0,s_0,s_0,r,r)$ and
$(0,a,r,r,D,D)$. Their distance is
$\max(r-s_0,D-r)=D-r$, since
$r-s_0=r(D-r)/(3D-r)<D-r$. The local inverse and local Fisher law of
Section~\ref{sec:cubic} are uniform on a compact segment of the base
path through zero. Below the entrance threshold all competitors are in
that base neighbourhood, proving \eqref{eq:B-below}. Above it, the
localization proof for the remote cone bounds its exact coordinates by
$O(\delta)$. The remote endpoint roots move by $O(\delta)$ and its
interior double cluster by $O(\sqrt\delta)$; the base inverse gives the
same $O(\sqrt\delta)$ root bound. A remote minimizer and the centre
are both feasible, and the exterior observation gap excludes every
other region. The bottleneck triangle inequality proves
\eqref{eq:B-above}.

For the last assertions, $b(s_0+\varepsilon)=D-\dot b\varepsilon+
O(\varepsilon^2)$. The function $R_s$ has a strict nondegenerate maximum
at $b(s)$ near the wall. Taylor expansion of $R_s(D)$ there, followed
by $x\mapsto x/(x-1)$, gives
\[
 c_D(s)-c_*(s)=
 \frac{-R_s''(b(s))}{2(R_s(b(s))-1)^2}(D-b(s))^2
                       +O((D-b(s))^3).
\]
This proves \eqref{eq:B-pencil-opening}. The strict weight margin keeps
$\gamma/\pi>c_D(s)$ for small $\varepsilon$, so the entire interval
$[c_*(s),c_D(s)]$ is feasible. Its upper endpoint has a root exactly at
$D$. A remote spectrum has increasing ordering
$(0,a_c,r,r,b_-(c),b_+(c))$, with $s<a_c<r<b_-(c)\le b_+(c)\le D$.
Its distance to the base is exactly
$\max(r-s,b_+(c)-r)\le D-r$, and equality is attained at $c_D(s)$.
The remote spectra converge uniformly to one spectrum as
$\varepsilon\downarrow0$; their mutual diameter is therefore less than
$D-r$ for small $\varepsilon$. Theorem~\ref{thm:cubic-fibre} excludes
any other exact component. This proves \eqref{eq:B-exact-diameter}.
\end{proof}
